\section{Summary, Convergence, and Outlook}

\subsection{Overview}

The preceding sections formulated structural constraints on intelligence-like dynamics. This final section summarizes those constraints, describes the convergence with independent research lines, and clarifies the scope of subsequent work.

\subsection{Summary of constraints}

\textbf{Constraint 1: Intelligence is not a stopping structure.}
Intelligence is not an ability, possession, or stored body of knowledge. It is a dynamical state that persists only under structural conditions.

\textbf{Constraint 2: Intelligence is not preserved, only regenerated.}
What can be stored are stopping products: knowledge, records, logs, and stabilized representations. Intelligence itself must be regenerated through ongoing dynamics.

\textbf{Constraint 3: Intelligence requires three simultaneous axiomatic conditions.}
Connected state space, temporal ordering, and global constraint $\Phi$ are jointly necessary for the minimal dynamics described here.

\textbf{Constraint 4: Correctness does not stabilize intelligence.}
Correctness increases reuse, reuse increases connectivity, and connectivity produces branching explosion.

\textbf{Constraint 5: Individual closure is structurally unavailable.}
The individual inferential system does not contain its own stopping operator. Apparent stopping is interruption unless it is stabilized through externalized structure.

\textbf{Constraint 6: Crossing the local horizon-form is structurally necessary.}
The differential horizon of intelligence is crossed through log, response, delay, agreement, and distributed responsibility. These do not guarantee correctness; they permit stopping, persistence, and later revision.

These constraints are necessary conditions, not sufficient conditions. A system that satisfies them is not automatically highly intelligent. A theory that ignores them will tend to reproduce the instabilities described in this paper.

\subsection{The convergence of independent research lines}

The structural argument of this paper converges with several independent research paths.

The Physarum research lineage begins from biophysical observation. It describes adaptive behavior through flow reinforcement, decay, and global boundary conditions without attributing intelligence as a possessed property.

Transformer architectures begin from engineering optimization. Their attention mechanisms and layer structure nevertheless instantiate connected state space, ordered updates, and global training constraints.

The present paper begins from structural analysis. The convergence of these paths does not prove a final theory. It suggests that the same minimal constraints recur across substrates.

\subsection{Outlook}

This paper is not the completion of a theory of intelligence. It is the first part of an intelligence-layer analysis within the VED/IFGT series.

Subsequent work must address:

\begin{itemize}
\item externalized stopping mechanisms;
\item log-based persistence;
\item distributed revisability;
\item social and institutional structures;
\item collective decision-making;
\item AI safety problems, including paperclip-type runaway optimization;
\item the transition dynamics among biological, consciousness, intelligence, and social layers.
\end{itemize}

The consciousness paper is complementary. It treats temporal non-closure and log referencing inside the individual window. The present paper treats the point at which individual inference must be externalized. Bio-IFGT supplies the biological quasi-closure background, while VED Vol.4 supplies the parent principle of the \DH{}.

\subsection{Concluding remarks}

This paper is a structural paper. It offers minimal axioms, structural constraints, and a naming of the local horizon-form that appears at the intelligence layer. It is not a complete definition of intelligence, not an implementation protocol, not a normative theory, and not a replacement for cognitive science or AI research.

The minimal statement is:

\begin{quote}
Intelligence is neither ability nor possession nor accumulated correctness. It is a dynamical state that persists only under structural conditions, and those conditions cannot close within the individual system alone. Intelligence therefore carries with it a local horizon-form: the differential horizon of intelligence.
\end{quote}

In the broader series, VED Vols.1--4 formulate the physical differential-development structure, Bio-IFGT formulates biological quasi-closure, the consciousness paper formulates log referencing at the consciousness layer, and the present paper formulates the minimal structure and local horizon-form of intelligence. Part II will address the social layer in which externalized stopping mechanisms are constructed.
